\section{Experiments}
\label{sec:experiments}

Given the setup in Section~\ref{sec:setup}, three experiments answer three
separate questions, each tested against a specific, falsifiable claim.

\textbf{R1} asks whether Farid et al.'s baselines and hybrids reproduce
under the paper's own components. It tests two claims: that the C4.5 and
NB baselines replicate closely (Section~\ref{sec:results} reports the
per-dataset gap to the paper's Tables 8--11), and that protocol A inflates
Algorithm~1 more than Algorithm~2, since Algorithm~1's filter removes rows
from the training set that a leaky protocol also removes from the test
folds, while Algorithm~2 only drops columns, which changes what a fold
learns from but not which instances a fold ever sees.

\textbf{E3a} asks whether the R1 inflation pattern is specific to the
Weka/J48 implementation or a property of the algorithms and the leaky
protocol in general, by repeating the same protocol-A-vs-B comparison with
an independent sklearn/one-hot/CART implementation. If both
implementations show the same asymmetry (large Algorithm~1 inflation,
small Algorithm~2 inflation), the asymmetry is attributable to the
algorithms' design, not to one toolchain.

\textbf{E9} asks two questions about the pipeline that Farid et al.\ never
built: whether either honest hybrid beats its non-hybrid baseline under
protocol B (Algorithm~1's tree vs.\ a plain tree; Algorithm~2's NB vs.\
plain NB), and whether the order in which the two algorithms are chained
(Path~1, filter-then-weight, vs.\ Path~2, weight-then-filter) changes
accuracy. E9 also records, per dataset, the percentage of training rows
Algorithm~1 removes and the percentage of attributes Algorithm~2 keeps, to
check whether a large removal or reduction predicts a large accuracy
change in the downstream classifier, which is the relationship plotted in
Fig.~\ref{fig:removal}.

All three experiments share the same ten datasets, the same protocol-B
refit-per-fold discipline as the default honest condition, and the same
Wilcoxon-over-datasets testing procedure described in
Section~\ref{sec:setup}. Table~\ref{tab:tests} collects every statistical
test from all three experiments in one place so the eight $p$-values can
be Holm-adjusted as one family rather than read as eight independent
tests.
